\chapter{Supplementary Methodological Detail}
\label{app:supplementary-methods}

This appendix provides additional context for the registration, model configuration and diagnostic assessment reported in Chapter~\ref{chap:methods}.

\section{Registration and spatial quality control}

Figure~\ref{fig:supplementary-qc} shows synthetic registration examples. Each panel illustrates the kinds of scale and orientation variation considered during review and provides visual context for the registration procedure described in Section~\ref{sec:registration}.

\section{Model configuration and validation}

The primary and secondary analyses used separate model ladders because their OCT predictor representations were different. The primary analysis used image-derived representations for the full eligible cohort, whereas the secondary analysis used source-mapped boundary features in the restricted boundary-available cohort. Both analyses retained participant-grouped validation and training-only preprocessing, as described in Sections~\ref{sec:models} and~\ref{sec:validation}.

\section{Extended diagnostic assessment}

In addition to the headline performance measures, the analysis considered calibration, residual distributions, fold-level error and participant-level error. These diagnostic assessments provide context for the model comparisons in Chapter~\ref{chap:results}, particularly as the observations within a participant are correlated. B0--B4 remain neutral source-mapped labels throughout the secondary analysis because their anatomical provenance was not independently verified.

\begin{figure}[htbp]
  \centering
  \includegraphics[width=0.92\textwidth]{figure_supplementary_qc_public.png}
  \caption{Supplementary registration schematics using synthetic coordinates and abstract geometric curves.}
  \label{fig:supplementary-qc}
\end{figure}
